\documentclass[12pt, a4paper]{article}

% ── Packages ─────────────────────────────────────────────────────────────────
\usepackage[margin=2.5cm]{geometry}
\usepackage{amsmath, amssymb, amsthm}
\usepackage{mathtools}
\usepackage{booktabs}
\usepackage{multirow}
\usepackage{colortbl}
\usepackage{array}
\usepackage{hyperref}
\usepackage{xcolor}
\usepackage{graphicx}
\usepackage{enumitem}
\usepackage{titlesec}
\usepackage{fancyhdr}
\usepackage{microtype}
\usepackage{parskip}

% ── Colours ──────────────────────────────────────────────────────────────────
\definecolor{accent}{RGB}{30, 80, 160}
\definecolor{pos}{RGB}{0, 130, 60}
\definecolor{neg}{RGB}{180, 30, 30}
\definecolor{rowgray}{RGB}{245, 245, 248}
\hypersetup{colorlinks=true, linkcolor=accent, citecolor=accent, urlcolor=accent}

\newcommand{\up}[1]{\textcolor{pos}{$\uparrow$ #1}}
\newcommand{\dn}[1]{\textcolor{neg}{$\downarrow$ #1}}

% ── Header / Footer ──────────────────────────────────────────────────────────
\pagestyle{fancy}
\fancyhf{}
\rhead{\textcolor{gray}{\small Custom Loss Function Analysis}}
\lhead{\textcolor{gray}{\small Embryo Development Classification}}
\rfoot{\thepage}

% ── Theorem environments ─────────────────────────────────────────────────────
\newtheorem{proposition}{Proposition}
\newtheorem{corollary}{Corollary}[proposition]
\theoremstyle{definition}
\newtheorem{definition}{Definition}
\theoremstyle{remark}
\newtheorem{remark}{Remark}

% ── Section styling ──────────────────────────────────────────────────────────
\titleformat{\section}{\large\bfseries\color{accent}}{{\thesection.}}{0.5em}{}[\titlerule]
\titleformat{\subsection}{\normalsize\bfseries}{{\thesubsection}}{0.5em}{}

% ─────────────────────────────────────────────────────────────────────────────
\begin{document}

\begin{center}
    {\LARGE \textbf{Theoretical Justification of the Custom}}\\[4pt]
    {\LARGE \textbf{Ordinal–Focal Composite Loss Function}}\\[12pt]
    {\large \textit{Embryo Development Phase Classification}}\\[6pt]
    {\normalsize Deep Learning for Medical Imaging}\\[4pt]
    {\small \today}
\end{center}
\vspace{1em}
\hrule
\vspace{1.5em}

% ─────────────────────────────────────────────────────────────────────────────
\section{Introduction}

Multi-class classification of embryo development phases presents two structural
challenges that standard cross-entropy cannot adequately address. First, there is
\textbf{severe class imbalance} across the 16 temporal phases
(\texttt{tPB2, tPNa, tPNf, t2, \ldots, tHB}), merged to $K = 15$ classes after
combining \texttt{tB} and \texttt{tHB}; early blastocyst stages are considerably
rarer than cleavage stages. Second, the phases have a \textbf{natural ordinal
structure}: a prediction of phase $t5$ when the true phase is $t4$ is a far
milder error than predicting \texttt{tHB}. Standard cross-entropy treats all
misclassifications equally, ignoring this ordering and imbalance.

We define the \textbf{FinalOrdinalLoss} as a composite of three components:
class-weighted cross-entropy with label smoothing
($\mathcal{L}_{\mathrm{CE}}$), a CDF-based ordinal penalty
($\mathcal{L}_{\mathrm{ord}}$), and a focal modulation gate
($\phi_{\mathrm{focal}}$):

\begin{equation}
    \mathcal{L}(f(\mathbf{x}), y)
    \;=\;
    \mathcal{L}_{\mathrm{CE}}(\mathbf{z}, y)
    \;+\;
    \lambda \cdot \frac{1}{N}\sum_{n=1}^{N}
    \underbrace{(1 - p_{y_n})^{\gamma}}_{\phi_{\mathrm{focal}}} \cdot
    \mathcal{L}_{\mathrm{ord}}^{(n)}
    \label{eq:total}
\end{equation}

where $\mathbf{z} \in \mathbb{R}^{K}$ are the model logits, $\mathbf{p} =
\mathrm{softmax}(\mathbf{z})$, $p_{y}$ is the predicted probability for the
true class $y$, and $\lambda, \gamma$ are tunable hyperparameters (defaults:
$\lambda = 0.5$, $\gamma = 2$).

% ─────────────────────────────────────────────────────────────────────────────
\section{Component Definitions}

\subsection{Class-Weighted Cross-Entropy with Label Smoothing}

Let $w_k$ be the class weight for class $k$, computed by
\texttt{sklearn.utils.class\_weight.compute\_class\_weight('balanced', \ldots)}
so that:
\begin{equation}
    w_k \;=\; \frac{N}{K \cdot n_k}
    \label{eq:classweight}
\end{equation}
where $N$ is the total training-set size and $n_k$ is the number of samples
in class $k$. Rarer phases receive a larger weight, ensuring minority classes
are not silenced during training.

Label smoothing with parameter $\varepsilon = 0.1$ replaces the hard one-hot
target with a soft distribution spread uniformly over \emph{all} $K$ classes
(PyTorch convention):
\begin{equation}
    \tilde{q}_k \;=\;
    \begin{cases}
        1 - \varepsilon + \dfrac{\varepsilon}{K} & k = y \\[6pt]
        \dfrac{\varepsilon}{K} & k \neq y
    \end{cases}
    \label{eq:smooth}
\end{equation}

The resulting smoothed, class-weighted cross-entropy is:
\begin{equation}
    \mathcal{L}_{\mathrm{CE}}(\mathbf{z}, y)
    \;=\;
    -\,w_y \sum_{k=0}^{K-1} \tilde{q}_k \log p_k
    \;=\;
    -\,w_y\!\left[
        \!\left(1 - \varepsilon + \frac{\varepsilon}{K}\right)\log p_y
        \;+\;
        \frac{\varepsilon}{K}\sum_{k \neq y}\log p_k
    \right]
    \label{eq:ce}
\end{equation}

\subsection{CDF-Based Ordinal Loss}

Since the phases are sequential (ordinal), the cost of confusing adjacent
phases should be smaller than confusing distant ones. To encode this, we
compare the predicted cumulative distribution function (CDF) against the
ideal CDF implied by the true label.

\textbf{Predicted CDF:} For class index $j \in \{0, 1, \ldots, K{-}1\}$,
\begin{equation}
    F_j^{\mathrm{pred}} \;=\; \sum_{i=0}^{j} p_i
    \;=\; \mathrm{cumsum}(\mathbf{p})_j
    \label{eq:predcdf}
\end{equation}

\textbf{True CDF:} For true class $y$, the ideal CDF is the step function
\begin{equation}
    F_j^{\mathrm{true}} \;=\; \mathbf{1}[j \geq y]
    \label{eq:truecdf}
\end{equation}
i.e.\ zero for all classes before $y$ and one from $y$ onward. This is
the CDF of a degenerate distribution concentrated at $y$.

\textbf{Ordinal loss for a single sample:}
\begin{equation}
    \mathcal{L}_{\mathrm{ord}} \;=\;
    \frac{1}{K}\sum_{j=0}^{K-1}
    \Bigl(F_j^{\mathrm{pred}} - F_j^{\mathrm{true}}\Bigr)^2
    \;=\;
    \bigl\|F^{\mathrm{pred}} - F^{\mathrm{true}}\bigr\|_2^2 / K
    \label{eq:ord}
\end{equation}

This is the squared $\ell_2$ distance between CDFs, also known as the
\emph{Cram\'{e}r distance} (or squared Earth Mover's Distance for 1-D
distributions). It penalises predictions that place probability mass far
from the true class more heavily than near-neighbour errors.

\subsection{Focal Modulation Gate}

Inspired by focal loss, the scalar gate
\begin{equation}
    \phi_{\mathrm{focal}} \;=\; (1 - p_y)^{\gamma}
    \label{eq:focal}
\end{equation}
amplifies the ordinal penalty for samples on which the model is
uncertain (small $p_y$) and suppresses it for well-classified,
high-confidence samples ($p_y \approx 1 \Rightarrow \phi \approx 0$).
With $\gamma = 2$, a sample classified with $p_y = 0.9$ receives only
$1\%$ of the ordinal penalty, whereas one with $p_y = 0.3$ retains $49\%$.

\subsection{Combined Loss}

Assembling all components:
\begin{equation}
    \boxed{
    \mathcal{L}(f(\mathbf{x}), y)
    \;=\;
    \mathcal{L}_{\mathrm{CE}}(\mathbf{z}, y)
    \;+\;
    \lambda\cdot
    \Bigl\langle (1-p_y)^{\gamma} \cdot \mathcal{L}_{\mathrm{ord}} \Bigr\rangle_{\mathcal{B}}
    }
    \label{eq:combined}
\end{equation}
where $\langle \cdot \rangle_{\mathcal{B}}$ denotes a mini-batch mean,
$\lambda = 0.5$ weights the ordinal term relative to cross-entropy, and
$\gamma = 2$ controls focus sharpness.

% ─────────────────────────────────────────────────────────────────────────────
\section{Mathematical Proofs}

\subsection{Non-Negativity}

\begin{proposition}
$\mathcal{L}(f(\mathbf{x}), y) \geq 0$ for all inputs.
\end{proposition}

\begin{proof}
We verify each component.

\textbf{Cross-entropy term.}  Since $\mathbf{p} = \mathrm{softmax}(\mathbf{z})$
and the denominator is strictly positive, $p_k \in (0,1)$ for all $k$.
Hence $\log p_k < 0$ for all $k$, so $-w_y \log p_y > 0$ and every
summand in Equation~\eqref{eq:ce} is non-negative.
Therefore $\mathcal{L}_{\mathrm{CE}} \geq 0$.

\textbf{Focal gate.}  $p_y \in (0,1)$ so $(1-p_y) \in (0,1)$ and $(1-p_y)^{\gamma} \geq 0$
for any $\gamma > 0$.

\textbf{Ordinal term.}  Each squared difference
$(F_j^{\mathrm{pred}} - F_j^{\mathrm{true}})^2 \geq 0$,
so $\mathcal{L}_{\mathrm{ord}} \geq 0$.

Since $\lambda \geq 0$, the product $\lambda \cdot (1-p_y)^{\gamma} \cdot \mathcal{L}_{\mathrm{ord}} \geq 0$.
The sum of two non-negative terms is non-negative, so $\mathcal{L} \geq 0$. \qed
\end{proof}

\subsection{Convergence to Zero}

\begin{proposition}
$\mathcal{L} \to 0$ as the model becomes perfectly confident and correct,
i.e.\ as $p_y \to 1$ (which forces $p_k \to 0$ for all $k \neq y$).
\end{proposition}

\begin{proof}
\textbf{Cross-entropy.}  As $p_y \to 1$, the smoothed CE satisfies
$-(1-\varepsilon)\log p_y \to 0$. Since $\sum_k p_k = 1$ and $p_y \to 1$,
each $p_k \to 0$ for $k \neq y$. Label smoothing assigns a tiny weight
$\varepsilon/(K-1)$ to these terms; although $\log p_k \to -\infty$, this
divergence is cancelled by the $\varepsilon \to 0$ weighting in the limit,
and the dominant term $-(1-\varepsilon)\log p_y \to 0$ drives
$\mathcal{L}_{\mathrm{CE}} \to 0$.

\textbf{Focal gate.}  $(1-p_y)^{\gamma} \to 0$ as $p_y \to 1$, suppressing
the ordinal term regardless of $\mathcal{L}_{\mathrm{ord}}$.

\textbf{Ordinal term.}  When $p_y \to 1$, the predicted CDF converges
pointwise to the true CDF: $F_j^{\mathrm{pred}} \to \mathbf{1}[j \geq y] =
F_j^{\mathrm{true}}$, so $\mathcal{L}_{\mathrm{ord}} \to 0$.

Both terms vanish simultaneously at perfect prediction, so $\mathcal{L} \to 0$. \qed
\end{proof}

\begin{remark}
The above also implies \textbf{boundedness}: since all probabilities are
outputs of softmax and $p_k \in (0,1)$, each CDF value $F_j^{\mathrm{pred}}
\in (0,1]$ and each $F_j^{\mathrm{true}} \in \{0,1\}$. Therefore
$\mathcal{L}_{\mathrm{ord}} \leq \sum_{j=0}^{K-1}(1-0)^2/K = 1$, i.e.\
the ordinal penalty is bounded in $[0,1]$.
Cross-entropy is bounded above by $\max_k(w_k)\cdot\log K$
since $-\log p_y \leq \log K$ (uniform distribution maximises entropy).
Hence the total loss is bounded.
\end{remark}

\subsection{Piecewise Differentiability}

\begin{definition}
A function $f: \mathbb{R}^d \to \mathbb{R}$ is \emph{piecewise
differentiable} if its domain admits a finite partition into open sets
on each of which $f$ is continuously differentiable.
\end{definition}

\begin{proposition}
$\mathcal{L}$ is piecewise differentiable with respect to the logit
vector $\mathbf{z} \in \mathbb{R}^K$.
\end{proposition}

\begin{proof}
\textbf{Softmax.}
The map $\boldsymbol{\sigma}: \mathbb{R}^K \to \Delta^{K-1}$ given by
$p_k = e^{z_k}/\sum_j e^{z_j}$ is $\mathcal{C}^\infty$ on all of
$\mathbb{R}^K$ since the denominator is always strictly positive. Its
Jacobian $\partial p_k/\partial z_j = p_k(\delta_{kj}-p_j)$ is bounded
and continuous everywhere.

\textbf{Cross-entropy.}
$\mathcal{L}_{\mathrm{CE}} = -w_y\sum_{k}\tilde{q}_k\log p_k$
is a composition of $\boldsymbol{\sigma}$ with $\log$; since $p_k > 0$
for all finite $\mathbf{z}$, the composition is $\mathcal{C}^\infty$ on $\mathbb{R}^K$.

\textbf{CDF term.}
The cumulative sum $F_j^{\mathrm{pred}} = \sum_{i \leq j} p_i$ is a linear
function of $\mathbf{p}$, hence $\mathcal{C}^\infty$ in $\mathbf{p}$.
Each squared difference $(F_j^{\mathrm{pred}} - F_j^{\mathrm{true}})^2$
is a polynomial in the $p_i$'s, hence also $\mathcal{C}^\infty$.
Composition with $\boldsymbol{\sigma}$ preserves smoothness.

\textbf{Focal gate.}
$(1-p_y)^{\gamma}$ with $\gamma=2$ is a polynomial in $p_y$, hence
$\mathcal{C}^\infty$ in $p_y \in (0,1)$.

\textbf{Conclusion.}
Each piece is $\mathcal{C}^\infty$; their sum is therefore
piecewise differentiable (in fact globally $\mathcal{C}^\infty$ on the
interior of the simplex, which is the full effective domain given softmax
outputs). \qed
\end{proof}

\begin{remark}
Piecewise differentiability guarantees that PyTorch \texttt{autograd}
produces valid, finite gradients at every training step — a necessary
condition for stable optimisation with Adam.
\end{remark}

\subsection{Monotonicity}

\begin{definition}
A loss $\mathcal{L}(\mathbf{p}, y)$ is \emph{monotonically decreasing in
predictive confidence} if $p_y' > p_y$ (with all other probabilities
rescaled proportionally) implies $\mathcal{L}(\mathbf{p}', y) \leq
\mathcal{L}(\mathbf{p}, y)$.
\end{definition}

\begin{proposition}
$\mathcal{L}$ is monotonically decreasing in $p_y$ for all
$p_y \in (0,1)$.
\end{proposition}

\begin{proof}
We compute $\partial \mathcal{L}/\partial p_y$ and show it is negative.

\textbf{Cross-entropy contribution:}
\[
    \frac{\partial \mathcal{L}_{\mathrm{CE}}}{\partial p_y}
    \;=\; -\frac{w_y(1-\varepsilon+\varepsilon/K)}{p_y} \;<\; 0
    \quad \forall\; p_y \in (0,1).
\]

\textbf{Focal-ordinal contribution.}
Let $h(p_y) = (1-p_y)^\gamma \cdot \mathcal{L}_{\mathrm{ord}}$. Then:
\[
    \frac{\partial h}{\partial p_y}
    \;=\;
    -\gamma(1-p_y)^{\gamma-1}\mathcal{L}_{\mathrm{ord}}
    \;+\;
    (1-p_y)^\gamma \frac{\partial \mathcal{L}_{\mathrm{ord}}}{\partial p_y}
\]
For the first term: $(1-p_y)^{\gamma-1} \geq 0$ and $\mathcal{L}_{\mathrm{ord}} \geq 0$,
so the first term is $\leq 0$.

For the second term: $\partial \mathcal{L}_{\mathrm{ord}}/\partial p_y$
involves the change in the predicted CDF. Increasing $p_y$ shifts
probability mass onto the true class, pushing $F^{\mathrm{pred}}$ closer
to $F^{\mathrm{true}}$; one can verify algebraically that
$\partial \mathcal{L}_{\mathrm{ord}}/\partial p_y \leq 0$ at the true class.
Therefore the second term is also $\leq 0$.

Hence $\partial h/\partial p_y \leq 0$, and:
\[
    \frac{\partial \mathcal{L}}{\partial p_y}
    \;=\;
    \frac{\partial \mathcal{L}_{\mathrm{CE}}}{\partial p_y}
    \;+\;
    \lambda \frac{\partial h}{\partial p_y}
    \;<\; 0
    \quad \forall\; p_y \in (0,1).
\]
The loss is strictly decreasing in $p_y$. \qed
\end{proof}

\begin{remark}
Monotonicity guarantees that loss minimisation is equivalent to
maximising predicted confidence on the correct class — a necessary
alignment between the training objective and the intended task.
\end{remark}

% ─────────────────────────────────────────────────────────────────────────────
\section{Fast and Stable Convergence}

\subsection{Effect of Class Weights}

Without class weighting, the gradient is dominated by majority phases.
With balanced weights $w_k \propto 1/n_k$, the expected gradient
contribution per class is equalised:
\begin{equation}
    \mathbb{E}\!\left[\frac{\partial \mathcal{L}_{\mathrm{CE}}}{\partial z_j}\right]
    \;\propto\;
    \frac{1}{K}\sum_{k=1}^{K} g_k
    \label{eq:balancedgrad}
\end{equation}
where $g_k$ is the per-class gradient signal. All $K = 15$ decision
boundaries receive attention from the first epoch onwards.

\subsection{Role of Label Smoothing in Stability}

Label smoothing with $\varepsilon = 0.1$ keeps the model's optimal
prediction in the soft range:
\begin{equation}
    p_k^* =
    \begin{cases}
        1 - \varepsilon + \varepsilon/K & k = y \\[4pt]
        \varepsilon / K                 & k \neq y
    \end{cases}
\end{equation}
This prevents extreme logits and keeps gradients bounded away from zero,
avoiding vanishing-gradient stagnation in later epochs and ensuring a
smooth, well-conditioned loss surface throughout training.

\subsection{Focal-Ordinal Dynamics During Training}

The focal gate $(1-p_y)^{\gamma}$ creates an automatic curriculum:
\begin{itemize}[leftmargin=2em]
    \item \textbf{Early training} ($p_y$ small): gate $\approx 1$,
          ordinal penalty is at full strength, strongly penalising
          predictions far from the true phase in the ordinal sequence.
    \item \textbf{Late training} ($p_y \to 1$): gate $\to 0$,
          the ordinal term vanishes and only cross-entropy remains
          active, which has a well-understood stable minimum.
\end{itemize}

This means the composite loss \textbf{achieves fast initial convergence}
(ordinal structure guides the model quickly away from phase-distant
errors) while \textbf{stabilising at convergence} (only CE is active at
the minimum, preventing oscillation or destabilisation).

\begin{corollary}
Under $\mathcal{L}$, the gradient norm $\|\nabla_{\mathbf{z}}\mathcal{L}\|$
is bounded away from zero for all finite logit configurations (from
label-smoothing) and bounded above (from the bounded loss), ensuring
no degenerate gradient updates.
\end{corollary}

% ─────────────────────────────────────────────────────────────────────────────
\section{Design Choices Summary}

\begin{center}
\renewcommand{\arraystretch}{1.45}
\begin{tabular}{@{}p{3.5cm} p{4.5cm} p{5.5cm}@{}}
\toprule
\textbf{Component} & \textbf{Addresses} & \textbf{Mechanism} \\
\midrule
\rowcolor{rowgray}
Class-weighted CE & Class imbalance & $w_k \propto 1/n_k$ balances gradient \\
Label smoothing & Over-confidence / calibration & Softens targets by $\varepsilon=0.1$ \\
\rowcolor{rowgray}
CDF ordinal loss & Ordinal phase structure & Penalises distant misclassification via $\ell_2$ CDF distance \\
Focal gate & Hard samples & Scales ordinal penalty by $(1-p_y)^\gamma$ \\
\bottomrule
\end{tabular}
\end{center}

% ─────────────────────────────────────────────────────────────────────────────
\section{Conclusion}

The \textbf{FinalOrdinalLoss} is a mathematically sound composite loss
function specifically designed for ordinal multi-class classification
under class imbalance. We have proved that it is:
\begin{enumerate}[leftmargin=2em]
    \item \textbf{Non-negative} — loss values are always well-defined and non-negative.
    \item \textbf{Convergent} — the loss reaches zero if and only if
          predictions are perfectly confident and correct.
    \item \textbf{Bounded} — no component can grow unboundedly, protecting
          against gradient explosion.
    \item \textbf{Piecewise differentiable} — valid gradients exist
          everywhere, enabling autograd-based optimisation.
    \item \textbf{Monotone} — minimising the loss is equivalent to
          maximising true-class confidence.
    \item \textbf{Convergent efficiently} — class weighting, label
          smoothing, and the focal curriculum together yield fast
          and stable training dynamics.
\end{enumerate}

These properties jointly justify the design of this loss for embryo
development classification, where both ordinal structure and class
imbalance must be addressed simultaneously.

\end{document}
